% Example output — Dr. Jordan Chen applying to Whitfield University
% Generated by claude-resume-kit for demonstration purposes
% This is a fictional researcher; all data is fabricated.

\documentclass{resume}
\usepackage{hyperref}
\usepackage{enumitem}
\usepackage{fontawesome}
\usepackage{tikz}
\usepackage{graphicx}
\hypersetup{
    colorlinks = true,
    linkcolor = [rgb]{0.9,0.4,0.4},
    anchorcolor = [rgb]{0.9,0.4,0.4},
    citecolor = [rgb]{0.4,0.4,0.4},
    filecolor = [rgb]{0.4,0.4,0.4},
    urlcolor = [rgb]{0.0,0.0,0.99},
}
\usepackage{xcolor}
\usepackage[version=4,arrows=pgf-filled]{mhchem}
\usepackage[includefoot,left=0.5in,top=0.5in,right=0.5in,bottom=0.2in,textwidth=7.5in,textheight=10.8in]{geometry}
\usepackage{fancyhdr}
\pagestyle{fancy}
\fancyhf{}
\renewcommand{\headrulewidth}{0pt}
\fancyfoot[R]{\hfill \thepage/\pageref{LastPage}}
\newcommand{\tab}[1]{\hspace{.2667\textwidth}\rlap{#1}}
\newcommand{\itab}[1]{\hspace{0em}\rlap{#1}}

\name{Jordan Chen, Ph.D.}
\address{jordan.chen@email.com \\ +1 (555) 123-4567}
\address{Richland, WA (Open to relocation to Westbrook, MA)}
\address{{Computational Protein Engineering $\vert$ ML-Guided Enzyme Design $\vert$ Biomolecular Simulation}}

\begin{document}

\vspace{-0.15cm}

%----------------------------------------------------------------------------------------
%	SUMMARY
%----------------------------------------------------------------------------------------
\begin{rSection}{Summary}
Computational biologist with 8+ years combining \textbf{protein language models} and \textbf{molecular dynamics simulations} for enzyme engineering and drug discovery. Fine-tuned ESM-2 on 45K experimental stability measurements to screen 8,500 enzyme variants at 3,000$\times$ experimental throughput, with 5 hits confirmed by collaborators. Co-developed open-source transfer learning framework adopted by 4 external groups. 15 publications (7 first-author) in ACS Catalysis, J.\ Chem.\ Theory Comput., and J.\ Med.\ Chem.
\end{rSection}
\vspace{-0.15cm}

%----------------------------------------------------------------------------------------
%	TECHNICAL SKILLS
%----------------------------------------------------------------------------------------
\begin{rSection}{Technical Skills}

\begin{skillgroup}{Molecular Simulation \& Modeling}
\skilldash{\textbf{GROMACS}, OpenMM, AMBER -- metadynamics, replica exchange MD, free energy perturbation, umbrella sampling}
\skilldash{AlphaFold2, Rosetta, AutoDock Vina -- protein structure prediction, homology modeling, and molecular docking}
\skilldash{CHARMM36m, AMBER ff19SB, OPLS-AA/M -- force field selection and benchmarking for disordered proteins}
\skilldash{Collective variable design, enhanced sampling protocol development, convergence analysis, MM/PBSA free energy}
\end{skillgroup}

\begin{skillgroup}{Machine Learning \& Data Science}
\skilldash{\textbf{Protein language models} (ESM-2, 650M params), graph neural networks, transfer learning, active learning loops}
\skilldash{\textbf{PyTorch}, scikit-learn, BioPython -- model fine-tuning, embedding extraction, feature engineering, sequence analysis}
\skilldash{Regression, classification, cross-validation, Spearman/RMSE benchmarking, dataset curation from public databases}
\end{skillgroup}

\begin{skillgroup}{Programming \& HPC}
\skilldash{\textbf{Python}, Bash, SQL -- scientific computing, data pipelines, automated analysis workflows, database management}
\skilldash{\textbf{SLURM}, Snakemake, Git, DVC -- HPC job scheduling, workflow automation, version control, reproducible research}
\end{skillgroup}

\begin{skillgroup}{Analysis \& Visualization}
\skilldash{MDAnalysis, ProDy, PyMOL, matplotlib, seaborn -- trajectory analysis, structural visualization, publication figures}
\skilldash{PostgreSQL, pandas, NumPy -- curated stability databases with automated quality filters for ML pipelines}
\end{skillgroup}

\begin{skillgroup}{Domain Expertise}
\skilldash{Protein engineering, enzyme thermostability, folding thermodynamics, drug discovery, virtual screening workflows}
\skilldash{Intrinsically disordered proteins, ligand binding free energy, biocatalysis, directed evolution, rational design}
\end{skillgroup}

\end{rSection}
\vspace{-0.15cm}

%----------------------------------------------------------------------------------------
%	RESEARCH EXPERIENCE
%----------------------------------------------------------------------------------------
\begin{rSection}{Research Experience}

\begin{rSubsection}{ML-Accelerated Protein Engineering and Computational Enzyme Design}{\textcolor{black!60}{Aug 2023 -- Present}}{Postdoctoral Research Associate, Lakewood University}{}
\item Fine-tuned ESM-2 protein language model on 45K experimental melting temperatures, achieving 0.82 Spearman correlation and enabling 3,000$\times$ throughput screening of 8,500 enzyme variants for industrial thermostability.
\item Co-developed transfer learning framework from protein language models reducing labeled training data by 60\% across 5 enzyme families, released as open-source tool with 200+ GitHub stars.
\item Extended protein language model to predict enzyme solvent tolerance across 8 organic co-solvent systems, validating against 50-ns explicit-solvent MD for 80 enzyme variants and identifying 4 candidates for green chemistry.
\item Automated sequence-to-simulation pipeline using Snakemake workflow manager, reducing per-variant setup from 4 hours to 10 minutes and supporting 6 researchers across 3 active projects.
\item Revealed sequence-dependent enzyme unfolding pathway divergence at 340 K through 200-ns replica exchange MD simulations, identifying stabilizing salt bridge networks that informed rational design criteria.
\item Mentored 3 graduate students on protein ML pipelines and MD simulation workflows, with 1 student co-authoring a peer-reviewed publication within 8 months of joining.
\end{rSubsection}

\begin{rSubsection}{Enhanced Sampling Methods for Protein Folding and Ligand Binding}{\textcolor{black!60}{Aug 2018 -- Jul 2023}}{Ph.D.\ Researcher, Westfield Institute of Technology}{}
\item Developed metadynamics-based enhanced sampling protocol for protein folding free energy landscapes, predicting folding temperatures within 8 K of experiment across 6 small proteins.
\item Calculated relative binding free energies for 40 congeneric ligand pairs via free energy perturbation, achieving 0.9 kcal/mol RMSE against experimental IC50 data across 3 drug target families.
\item Built curated protein thermostability database integrating 12,000 experimental melting temperatures from 3 public sources, with automated quality filters adopted by 8 lab members for ML training set construction.
\item Benchmarked 4 protein force fields on 15 intrinsically disordered protein sequences, establishing CHARMM36m as the optimal choice for IDP conformational ensemble prediction with 40\% better agreement with SAXS data.
\end{rSubsection}

\begin{rSubsection}{Computational Biophysics and Structural Analysis}{\textcolor{black!60}{May 2016 -- Jul 2018}}{Undergraduate Research Assistant, Eastgate University}{}
\item Performed homology modeling and 100-ns MD simulations of 4 mutant lysozyme variants, identifying destabilizing cavity mutations consistent with published experimental unfolding data.
\item Built Python analysis scripts for automated hydrogen bond occupancy tracking across 500-ns aggregate trajectories, adopted by 3 lab members for ongoing protein stability projects.
\end{rSubsection}

\end{rSection}
\vspace{-0.15cm}

%----------------------------------------------------------------------------------------
%	EDUCATION
%----------------------------------------------------------------------------------------
\begin{rSection}{Education}
{Ph.D., Biomedical Engineering} \hfill {\textcolor{black!60}{Aug 2018 -- Jul 2023}}\\
{Westfield Institute of Technology}, Westfield, MA \hfill GPA: \textbf{3.92}/4.00

{B.S., Biochemistry (Honors)} \hfill {\textcolor{black!60}{Aug 2014 -- May 2018}}\\
{Eastgate University}, Portland, OR \hfill GPA: \textbf{3.87}/4.00
\end{rSection}
\vspace{-0.15cm}

%----------------------------------------------------------------------------------------
%	SELECTED PUBLICATIONS
%----------------------------------------------------------------------------------------
\begin{rSection2}{Selected Publications (15 papers $\vert$ 280+ citations)}

\item \textbf{J.\ Chen}, R.\ Nakamura, S.\ Patel, K.\ Holmberg, M.\ Rivera. ``Deep Learning-Guided Screening of Thermostable Enzyme Variants for Industrial Biocatalysis.'' \textit{ACS Catalysis}, 2025.

\item \textbf{J.\ Chen}, M.\ Rivera, K.\ Holmberg. ``Transfer Learning from Protein Language Models for Low-Data Enzyme Property Prediction.'' \textit{Bioinformatics}, 2024.

\item \textbf{J.\ Chen}, L.\ Alvarez. ``Ligand Binding Free Energies via Enhanced-Sampling FEP for Three Drug Target Families.'' \textit{J.\ Med.\ Chem.}, 2023.

\item \textbf{J.\ Chen}, L.\ Alvarez. ``Metadynamics Protocol for Protein Folding Free Energy Landscapes.'' \textit{J.\ Chem.\ Theory Comput.}, 2022.

\item \textbf{J.\ Chen}, P.\ Kowalski, L.\ Alvarez. ``Force Field Benchmarking for Intrinsically Disordered Protein Ensembles.'' \textit{J.\ Chem.\ Theory Comput.}, 2021.

\item \textbf{J.\ Chen}, P.\ Kowalski, L.\ Alvarez. ``Curated Thermostability Database for ML-Ready Protein Engineering Benchmarks.'' \textit{Bioinformatics}, 2021.

\item \textbf{J.\ Chen}, T.\ Yamamoto, K.\ Holmberg. ``Predicting Enzyme Solvent Tolerance with Fine-Tuned Protein Language Models.'' \textit{Proteins: Struct., Funct., Bioinf.}, 2025 (under review).

\end{rSection2}
\vspace{-0.15cm}

%----------------------------------------------------------------------------------------
%	HONORS & AWARDS
%----------------------------------------------------------------------------------------
\begin{rSection2}{Honors \& Awards}
\item \textbf{NSF Graduate Research Fellowship}, National Science Foundation (2019). Three-year fellowship supporting doctoral research in computational protein engineering.
\item \textbf{Best Oral Presentation}, Westfield Biophysics Symposium (2022). Enhanced sampling methods for protein folding thermodynamics.
\item \textbf{Dean's Teaching Award}, Westfield Institute of Technology (2021). Outstanding TA in computational biology.
\item \textbf{Outstanding Poster Award}, Gordon Research Conference on Proteins (2023). ML-guided enzyme thermostability screening.
\end{rSection2}
\vspace{-0.15cm}

%----------------------------------------------------------------------------------------
%	SELECTED PRESENTATIONS
%----------------------------------------------------------------------------------------
\begin{rSection2}{Selected Presentations}
\item ``ML-Guided Screening of Thermostable Enzyme Variants.'' \textit{Gordon Research Conference on Proteins}, Ventura, CA (2023). Poster.
\item ``Transfer Learning from Protein Language Models for Low-Data Enzyme Property Prediction.'' \textit{ACS National Meeting}, San Francisco, CA (2024). Oral.
\item ``Enhanced Sampling Protocols for Protein Folding Free Energy Landscapes.'' \textit{Biophysical Society Annual Meeting}, San Diego, CA (2022). Oral.
\item ``Force Field Benchmarking for Intrinsically Disordered Protein Ensembles.'' \textit{AIChE Annual Meeting}, Boston, MA (2021). Poster.
\item ``Automated Screening Pipelines for ML-Guided Enzyme Engineering.'' \textit{Lakewood University Computational Biology Seminar Series} (2024). Invited talk.
\end{rSection2}
\vspace{-0.1cm}

\begin{center}
\vspace{0.15cm}
\textit{Authorized to work in the United States}
\end{center}

\end{document}
